\documentclass[conference]{IEEEtran}
\IEEEoverridecommandlockouts

\usepackage{cite}
\usepackage{amsmath,amssymb,amsfonts,amsthm}
\usepackage{algorithm}
\usepackage{algpseudocode}
\usepackage{graphicx}
\usepackage{textcomp}
\usepackage{xcolor}
\usepackage{booktabs}
\usepackage{hyperref}
\usepackage{pgfplots}
\usepackage{balance}
\pgfplotsset{compat=1.18}

% Relax justification to avoid stretched word-spacing in narrow columns
\tolerance=2000
\emergencystretch=1.5em

\def\BibTeX{{\rm B\kern-.05em{\sc i\kern-.025em b}\kern-.08em
    T\kern-.1667em\lower.7ex\hbox{E}\kern-.125emX}}

\newtheorem{theorem}{Theorem}

% ─── METADATA ───────────────────────────────────────────────────────────────
\begin{document}

\title{Atomic Ladder Execution: A Serializable Protocol for Multi-Agent
Proxy Bidding in Real-Time Auction Platforms}

\author{
  \IEEEauthorblockN{Asha Joseph, Pratham Reet, Nitin A. Rane,
    Reeshav Sinha, and Om Ji Rao}
  \IEEEauthorblockA{
    \textit{Department of Computer Science and Engineering} \\
    \textit{New Horizon College of Engineering} \\
    Bengaluru, India \\
    prathamreet@gmail.com
  }
}

\maketitle

% ─── ABSTRACT ───────────────────────────────────────────────────────────────
\begin{abstract}
Every major online auction platform offers proxy bidding, where a
participant registers a private ceiling price and the system places
incremental counter-bids on their behalf until that ceiling is reached.
Despite widespread deployment, no published work specifies how competing
proxy agents are resolved when a new manual bid arrives. In common
practice, platforms compute the equilibrium price in closed form and
record a single database row, discarding the intermediate bidding
sequence. While efficient, this approach sacrifices two properties that
platform operators depend on: a faithful per-increment audit trail
suitable for dispute resolution and fraud analytics, and correct
serialisation when manual bidders interleave with ongoing proxy
resolution. An alternative approach that recurses through the bid
service for each increment introduces nested-transaction hazards under
standard object-relational mappers and risks partial-commit corruption.

This paper introduces the Atomic Ladder Protocol, which writes every
increment of a proxy-bid ladder as a distinct bid row within a single
database transaction. A two-tier lock ordering -- auction row first,
then wallets sorted by user identifier -- guarantees deadlock freedom
across all concurrent writers. We state five desired properties, prove
three correctness theorems (log fidelity, Vickrey-equivalent pricing,
and deadlock freedom), and benchmark the protocol on PostgreSQL~15 with
Redis~7. At 100 concurrent bidders, the Redis Stream variant sustains
975~bids/s with a 95th-percentile latency of 8.1~ms, whereas direct
row locking degrades to 30~bids/s at 5{,}772~ms -- a
$\mathbf{32.5\times}$ throughput gain.
\end{abstract}

\begin{IEEEkeywords}
Concurrency control, online auctions, proxy bidding, real-time systems,
serializability, two-phase locking, Vickrey auction.
\end{IEEEkeywords}

% ─── I. INTRODUCTION ────────────────────────────────────────────────────────
\section{Introduction}

Consider an English auction in which several participants have each
registered a private ceiling price. When a fresh manual bid arrives, the
platform must decide how many of those ceilings to exercise and what the
resulting price should be. Let $M_i$ denote the ceiling of bidder~$i$
and $\delta$ the minimum bid increment. Standard auction
theory~\cite{Vickrey1961} establishes that the stable equilibrium price
under truthful proxy delegation is
\begin{equation}
  P^\star = \min\!\bigl(M_{(1)},\; M_{(2)} + \delta\bigr)
  \label{eq:vickrey}
\end{equation}
where $M_{(1)} \geq M_{(2)} \geq \cdots$ ranks the ceilings in
descending order. The highest-ceiling bidder wins, paying only one
increment above the runner-up -- an outcome strategically equivalent to
a second-price sealed-bid auction.

Because~(\ref{eq:vickrey}) is a closed-form expression, the
implementation path taken by most deployed platforms is straightforward:
evaluate the formula, write one bid row at~$P^\star$, and broadcast the
result. External observations of eBay, Catawiki, and LiveAuctioneers
confirm that users see only the final jump, not the intermediate
steps. This {\em single-jump} strategy is $O(\log N)$ in the pool size
and generates minimal database load, yet it comes at a cost that is
seldom discussed in the literature.

The first casualty is {\em log fidelity}. A user who opens the bid
history sees the price leap from, say, Rs.\,1{,}000 to Rs.\,1{,}800 in
a single row. The incremental narrative -- the back-and-forth that
would have unfolded between competing proxies -- is absent from the
database. Several platforms attempt to reconstruct this sequence on the
client side for display purposes, but doing so produces no durable audit
trail. Regulatory compliance, dispute adjudication, and behavioural
fraud engines~\cite{Trevathan2007,Ford2010} all require the granular log
to reside in the database itself.

The second casualty is {\em serialisation under concurrent manual
bids}. If a human bidder submits a bid while a proxy ladder is being
evaluated, the platform must decide whether to block the newcomer (high
tail latency) or abort the ladder and restart (livelock risk under load).
Neither option is satisfactory at scale.

A natural alternative is to recurse through the bid-placement service
once per increment, recording each rung individually. In practice, this
approach is fragile. Under widely used object-relational
mappers~\cite{Prisma2024}, nesting one interactive transaction inside
another is either silently flattened or causes outright deadlocks. Worse,
should the process fail at rung~7 of a 12-rung ladder, the first six
rungs are already committed and the auction is left in a state from which
no replay path is obvious. We observed both failure modes during
development.

We propose the {\em Atomic Ladder Protocol}, which reconciles fidelity
with atomicity. Every rung of the proxy ladder is persisted as its own
bid row, yet the entire sequence commits or rolls back as one
transaction. Three claims are made:

\begin{enumerate}
  \item \textbf{Protocol design.} A bounded-iteration procedure that
    acquires the auction lock, enumerates the active proxy pool, and walks
    the price upward one increment per step within a single
    {\tt BEGIN}/{\tt COMMIT} scope. Wallet locks are acquired in ascending
    user-identifier order to preclude circular waits.

  \item \textbf{Formal guarantees.} We prove three properties: log
    fidelity (a contiguous price sequence with no gaps), Vickrey
    equivalence (the final price
    satisfies~(\ref{eq:vickrey}) regardless of arrival order), and
    deadlock freedom (no two writers obeying the same lock order can form
    a wait-for cycle).

  \item \textbf{Empirical evaluation.} Benchmarks using k6 compare the
    atomic ladder against a direct row-locking baseline at concurrency
    levels of 1, 10, and 100 virtual users. We also report determinism
    (hash-based log comparison across repeated runs) and log-fidelity
    counts.
\end{enumerate}

The protocol is deployed in the open-source \textsc{AuctionHaus}
platform. The reference implementation spans roughly 180~lines of
TypeScript, packaged as a BullMQ worker job that fires after each manual
bid.

% ─── II. RELATED WORK ────────────────────────────────────────────────────────
\section{Related Work}

\subsection{Proxy Bidding Mechanisms}

The strategic theory of proxy bidding traces back to
Vickrey's~\cite{Vickrey1961} analysis of second-price sealed tenders. In
the internet era, Roth and Ockenfels~\cite{Roth2002} examined eBay's and
Amazon's proxy mechanisms and showed that truthful ceiling submission is a
weakly dominant strategy -- making proxy-mediated English auctions
strategically equivalent to Vickrey's sealed-bid game. Their work focuses
on bidder behaviour, however, not on how the platform resolves the proxy
pool internally. That implementation detail is the gap we address.

\subsection{Concurrency Control in Transaction Systems}

Row-level pessimistic locking is a well-established tool for
serialisation, covered comprehensively by Bernstein, Hadzilacos, and
Goodman~\cite{Bernstein1987}. The specific technique of acquiring locks
in a predetermined total order to eliminate deadlocks is attributed to
Gray and Reuter~\cite{Gray1992}, who prove that a consistent lock
sequence makes wait-for cycles impossible. Kleppmann~\cite{Kleppmann2017}
surveys how these classical techniques map onto modern database isolation
levels and warns that PostgreSQL's default {\tt READ COMMITTED} isolation
permits lost-update anomalies in read-modify-write patterns unless
explicit row locks are taken -- precisely the scenario our protocol
targets.

\subsection{Auction Architecture Reports}

Published architecture descriptions of production auction platforms are
scarce. The eBay architecture overview by Shoup and
Pritchett~\cite{Shoup2004} predates the current proxy bidding
implementation. Ockenfels, Reiley, and Sadrieh~\cite{Ockenfels2006}
survey bidder behaviour in online auctions using log dumps, yet they do
not document the resolution algorithm that produced those logs. Catawiki,
Saleroom, and LiveAuctioneers publish only marketing material. To the
best of our knowledge, no peer-reviewed paper details how a deployed
auction platform resolves concurrent proxy bids while retaining a
per-increment log.

% ─── III. SYSTEM MODEL ────────────────────────────────────────────────────────
\section{System Model and Problem Statement}

\subsection{Objects and Operations}

The platform exposes three durable objects:

\begin{itemize}
  \item \textbf{Auction} $A$ carries five fields: start price,
    current price, minimum increment, status, and end time.
    A single row per auction ID.

  \item \textbf{Wallet} $W_u$ per user $u$ with
    $\{\text{balance},\,\text{heldAmount}\}$. The available balance of $u$
    is $\text{balance}_u - \text{heldAmount}_u$.

  \item \textbf{AutoBid} $\alpha_{u,A} = (u, A, M_{u,A})$: user $u$ has
    delegated up to $M_{u,A}$ for auction $A$. At most one per
    $(u,A)$ pair.
\end{itemize}

The platform exposes two write operations on this state:

\begin{itemize}
  \item \textbf{PlaceBid}$(u, A, p)$: user $u$ manually bids amount $p$
    on auction $A$. Atomically validates $p \geq \text{currentPrice}_A +
    \delta_A$, holds $p$ in $W_u$, and updates
    $\text{currentPrice}_A \leftarrow p$.

  \item \textbf{ResolveAutoBids}$(A, u_{\text{trig}})$: triggered after
    every successful PlaceBid. Walks the auto-bid pool, choosing
    challengers in priority order, and emits one bid row per ladder rung.
    This is the operation the protocol governs.
\end{itemize}

\subsection{Concurrency Assumptions}

We assume:

\begin{enumerate}
  \item PostgreSQL default isolation level (\texttt{READ COMMITTED}).
  \item \texttt{SELECT \ldots FOR UPDATE} blocks concurrent updaters of
    the locked row.
  \item Multiple worker processes consume from the same BullMQ
    \texttt{auto-bid} queue; per-auction serialisation comes from the
    database lock, not the queue's group-key feature.
  \item BullMQ may re-deliver any job (worker crash, stalled-job
    recovery), so ResolveAutoBids must be idempotent.
\end{enumerate}

\subsection{Desired Properties}

For a manual bid that triggers a ladder resolution starting at price
$p_0$ and ending at $P^\star$ (as defined by
Equation~\ref{eq:vickrey} over the current auto-bid pool $\mathcal{P}$):

\begin{itemize}
  \item \textbf{P1 -- Log Fidelity.} The bid log contains, in commit
    order, exactly one row at each price
    $p_0 + k\delta$ for $k = 1, 2, \ldots, K$ where $p_0 + K\delta =
    P^\star$.

  \item \textbf{P2 -- Vickrey Equivalence.} The final price after the
    ladder commits equals $P^\star$, and the winning row is owned by
    the bidder with limit $M_{(1)}$.

  \item \textbf{P3 -- Serializability.} Any execution interleaved with
    other writers (PlaceBid, withdraw, settlement) is equivalent to
    some serial execution.

  \item \textbf{P4 -- Deadlock Freedom.} The protocol cannot deadlock
    against any other writer that obeys the global lock order.

  \item \textbf{P5 -- Bounded Work.} The number of rungs is bounded by
    $K \leq \lceil (M_{(1)} - p_0) / \delta \rceil + 1$.
\end{itemize}

% ─── 4. PROTOCOL ────────────────────────────────────────────────────────────
\section{The Atomic Ladder Protocol}

\subsection{Lock-Order Invariant}

We declare a single global lock order shared by every write operation in
the system:

\begin{quote}
\emph{Acquire the auction row \texttt{FOR UPDATE} first, then every
participating wallet \texttt{FOR UPDATE} in ascending \texttt{userId}
lexicographic order.}
\end{quote}

This invariant is shared by PlaceBid, Withdraw, BuyNow, ConfirmPayment,
EndAuction, and ResolveAutoBids. As long as every writer obeys it, the
classical theorem on lock ordering applies: no cycle in the
wait-for graph is possible (Property~P4).

\subsection{Pseudocode}

Algorithm~\ref{alg:ladder} describes the procedure. The procedure is
called \texttt{ResolveAutoBids}$(A, u_{\text{trig}})$ where $u_{\text{trig}}$
is the user whose manual bid triggered the resolution.

\begin{algorithm}[t]
\caption{ResolveAutoBids$(A, u_{\text{trig}})$}
\label{alg:ladder}
\begin{algorithmic}[1]
\State \textbf{begin transaction} (READ COMMITTED)
\State $\textsc{SELECT id FROM auctions WHERE id} = A \textsc{ FOR UPDATE}$
\State $\text{auction} \gets$ load $A$ with $\text{type, status, endTime, currentPrice}, \delta$
\If{$\text{auction} = \bot$ \textbf{or} $\text{auction.status} \neq \text{ACTIVE}$}
  \State \textbf{commit}; \Return $\emptyset$
\EndIf
\If{$\text{now}() > \text{auction.endTime}$} \Comment{Phase E4}
  \State \textbf{commit}; \Return $\emptyset$
\EndIf
\If{$\text{auction.type} \neq \text{ENGLISH}$}
  \State \textbf{commit}; \Return $\emptyset$
\EndIf
\State $\mathcal{P} \gets$ active auto-bids on $A$, excluding $u_{\text{trig}}$,
       sorted by $M_i$ desc then $\text{createdAt}$ asc
\If{$\mathcal{P} = \emptyset$} \textbf{commit}; \Return $\emptyset$
\EndIf
\State $U \gets$ ascending sort of $\{u_{\text{trig}}\} \cup \{u_i : \alpha_i \in \mathcal{P}\}$
\ForAll{$u \in U$}
  \State $\textsc{SELECT id FROM wallets WHERE userId} = u \textsc{ FOR UPDATE}$
\EndFor
\State $b_{\text{trig}} \gets$ most recent WINNING bid by $u_{\text{trig}}$ on $A$
\State $\text{cur} \gets \text{auction.currentPrice}$;
       $\text{winner} \gets b_{\text{trig}}$
\State $K_{\max} \gets \lceil (M_{(1)} - \text{cur})/\delta \rceil + 1$
\State $\text{steps} \gets$ empty list
\For{$k = 1$ to $K_{\max}$}
  \State $\text{next} \gets \text{cur} + \delta$
  \State $\mathcal{P} \gets \{\alpha_i \in \mathcal{P}: M_i \geq \text{next} \wedge u_i \neq \text{winner.bidder}\}$
  \If{$\mathcal{P} = \emptyset$} \textbf{break} \EndIf
  \State $\alpha^\ast \gets \arg\max_{\alpha \in \mathcal{P}} M_i$
  \State $W \gets$ wallet of $u^\ast$ (already locked above)
  \If{$W.\text{balance} - W.\text{heldAmount} < \text{next}$}
    \State deactivate $\alpha^\ast$;
           $\mathcal{P} \gets \mathcal{P} \setminus \{\alpha^\ast\}$;
           \textbf{continue}
  \EndIf
  \State outbid $\text{winner}$; release $\text{winner.amount}$ in winner's wallet
  \State insert new bid row $(u^\ast, A, \text{next}, \text{WINNING}, \text{isAuto}=\text{true})$
  \State decrement $W.\text{balance}$ by $\text{next}$;
         increment $W.\text{heldAmount}$ by $\text{next}$
  \State append the new row to $\text{steps}$
  \State $\text{cur} \gets \text{next}$; $\text{winner} \gets$ new row
\EndFor
\If{$|\text{steps}| > 0$}
  \State $\textsc{UPDATE auctions SET currentPrice} = \text{cur} \textsc{ WHERE id} = A$
\EndIf
\State \textbf{commit}; \Return $\text{steps}$
\end{algorithmic}
\end{algorithm}

\subsection{Implementation Notes}

The reference implementation is in TypeScript on Prisma\,5. The
\texttt{SELECT \ldots FOR UPDATE} statements use Prisma's
\texttt{\$queryRaw} since the generated client has no native syntax for
row locks. The outer transaction uses
\texttt{prisma.\$transaction(\textit{interactive callback})} which maps
to a single PostgreSQL \texttt{BEGIN/COMMIT} pair.

After the transaction commits, the worker emits a single
\texttt{bid:ladder} Socket.io event carrying the entire \texttt{steps}
array. This is a per-ladder event rather than one event per rung; clients
animate the rungs from the array (see Section~\ref{sec:eval}). The
single-emit strategy reduces $N$ socket round-trips to one for an
$N$-rung ladder, which we observe matters more than expected at high
fan-out (a hot auction with 200 subscribed sockets and a 10-rung ladder
moves from 2000 socket frames to 200).

% ─── 5. CORRECTNESS ─────────────────────────────────────────────────────────
\section{Correctness}

\subsection{Property P1: Log Fidelity}

\begin{theorem}[Log Fidelity]
Let $p_0$ be the auction current price after the trigger manual bid
commits, and let $P^\star$ be the final price after ResolveAutoBids
returns. Then for every $p \in \{p_0 + \delta, p_0 + 2\delta, \ldots,
P^\star\}$, the bid log contains exactly one row at price $p$, and the
rows are committed in price-ascending order.
\end{theorem}

\begin{proof}
The for-loop in Algorithm~\ref{alg:ladder} starts at $\text{cur} = p_0$
and, on each non-skipped iteration, inserts a row at
$\text{next} = \text{cur} + \delta$ and assigns
$\text{cur} \leftarrow \text{next}$. Therefore the inserted prices form a
contiguous sequence $p_0 + \delta, p_0 + 2\delta, \ldots, p_0 +
K\delta$ where $K$ is the number of successful iterations.

If an iteration skips (affordability check fails on $\alpha^\ast$),
$\text{cur}$ is not advanced and $\alpha^\ast$ is removed from $\mathcal{P}$.
The next iteration retries with a smaller pool; the value
$\text{next} = \text{cur} + \delta$ is unchanged. So a skip does not
create a gap in the price sequence.

Termination at $P^\star$: when $\mathcal{P}$ shrinks to empty (no
remaining $\alpha$ has $M_i \geq \text{next}$ \emph{and} is not the
current winner), the loop breaks. The current winner is the bidder
with $M_{(1)}$ (by induction -- each iteration replaces the winner with
the highest-$M$ challenger still in the pool); the price at which they
stand alone is $M_{(2)} + \delta$ (or $M_{(1)}$ if $M_{(1)} < M_{(2)} +
\delta$, when the loop's affordability/limit check terminates earlier).
Both cases match Equation~\ref{eq:vickrey}.

Commit order: all rows are inserted in the same transaction. Within a
PostgreSQL transaction, inserts are visible to subsequent statements in
the same transaction in the order they were issued; on commit, they
become visible to outside readers atomically but with insertion-order
identifiers (\texttt{ctid} ordering). Other transactions reading the bid
log after commit see the rows in price-ascending order.
\end{proof}

\subsection{Property P2: Vickrey Equivalence}

\begin{theorem}[Vickrey Equivalence]
For any number of concurrent active auto-bids $\mathcal{P}$ on auction
$A$, after a manual bid by $u_{\text{trig}}$ at price $p_0$, the final
auction price after ResolveAutoBids commits equals
$P^\star = \min(M_{(1)},\, M_{(2)} + \delta)$, where $M_{(1)}, M_{(2)}$
are the two largest limits among
$\mathcal{P} \cup \{p_0\}$ (treating the manual bid's effective limit as
$p_0$).
\end{theorem}

\begin{proof}
By induction on the number of iterations $k$. After iteration $k$, the
current winner has the $k$th-largest limit among
$\mathcal{P} \cup \{p_0\}$, and the price is $p_0 + k\delta$.

The loop terminates either when $\mathcal{P}$ has only the winner left,
or when the next-highest challenger has $M_i < \text{cur} + \delta$. Both
conditions match Equation~\ref{eq:vickrey}: in the first case the winner
is alone with the next-best limit consumed; in the second case the
winner has out-priced the next-best by enough that no further increment
is profitable.

The lock on the auction row ensures no concurrent PlaceBid can advance
$\text{currentPrice}$ during the loop. If a concurrent PlaceBid was
queued, it waits at the FOR UPDATE statement until our transaction
commits, then re-reads $\text{currentPrice}$ and triggers a fresh
ladder resolution -- a clean serial extension of the auction history.
\end{proof}

\subsection{Properties P3 and P4: Serializability and Deadlock Freedom}

\begin{theorem}[Serializability]
Any execution of ResolveAutoBids interleaved with other writers
(PlaceBid, Withdraw, EndAuction, ConfirmPayment) that obey the global
lock order is equivalent to a serial execution.
\end{theorem}

\begin{proof}
All writers acquire the auction row lock before any wallet lock. The
auction row is the only object touched by every writer, so it acts as
the universal serialisation point: the order in which transactions
acquire the auction lock determines the serial schedule. Wallet locks
acquired second cannot create a cycle because the auction lock
strictly precedes them.
\end{proof}

\begin{theorem}[Deadlock Freedom]
No deadlock is possible between any two writers in the system that obey
the global lock order.
\end{theorem}

\begin{proof}
Each transaction acquires locks in the same total order: first the
auction (single resource), then wallets in ascending \texttt{userId}.
By Gray and Reuter~\cite{Gray1992} \S7.4, lock ordering eliminates
wait-for cycles.
\end{proof}

\subsection{Property P5: Bounded Work}

The for-loop variable $k$ ranges over $[1, K_{\max}]$ where
$K_{\max} = \lceil (M_{(1)} - p_0)/\delta \rceil + 1$. Each iteration is
$O(|\mathcal{P}|)$ wallet/bid writes plus an $O(1)$ database round-trip.
For a typical English auction with $\delta = 1\%$ of $p_0$ and the
highest auto-bid at $2\times p_0$, $K_{\max} \approx 100$, well below
the operational ceiling we enforce ($K_{\max} \leq 10^5$).

% ─── VI. EVALUATION ──────────────────────────────────────────────────────────
\section{Evaluation}
\label{sec:eval}

\subsection{Experimental Setup}

Two bid-processing pipelines are compared on the same hardware
(Windows~11, Node.js~20, PostgreSQL~15, Redis~7) under identical
workloads generated by k6~\cite{k6}. Each run lasts 15~seconds against a
single English auction priced at Rs.\,1{,}000 with
$\delta =$~Rs.\,100. Wallet balances are set high enough
(Rs.\,100{,}000{,}000) to eliminate fund-exhaustion as a confound.

\begin{itemize}
  \item \textbf{Direct (FOR UPDATE).} The Express endpoint acquires the
    auction row lock inside a Prisma interactive transaction, validates
    the bid, and commits. Under contention, competing transactions queue
    in the PostgreSQL lock manager.

  \item \textbf{Sequencer (Redis Stream).} Bids are appended to a
    per-auction Redis Stream via {\tt XADD}. A lone consumer
    ({\tt XREADGROUP}) drains the stream and applies each bid to
    PostgreSQL serially. Because only one writer ever touches the
    database, the {\tt FOR UPDATE} lock is never contested.
\end{itemize}

\subsection{Throughput and Latency}

Table~\ref{tab:throughput} summarises the measured results.

\begin{table}[h]
\centering
\caption{Measured Throughput and 95th-Percentile Latency}
\label{tab:throughput}
\begin{tabular}{lrrrrr}
\toprule
\textbf{Implementation} & \textbf{C} & \textbf{Iters}
  & \textbf{Thr.} & \textbf{p95} \\
 & & (15\,s) & bids/s & ms \\
\midrule
Direct (FOR UPDATE)         & 1    & 53      & 3.5    & 236    \\
Sequencer (Redis Stream)    & 1    & 86      & 5.7    & 134    \\
\midrule
Direct (FOR UPDATE)         & 10   & 143     & 9.5    & 1{,}284 \\
Sequencer (Redis Stream)    & 10   & 1{,}392 & 92.8   & 48     \\
\midrule
Direct (FOR UPDATE)         & 100  & 570     & 30.0   & 5{,}772 \\
Sequencer (Redis Stream)    & 100  & 14{,}626& 975.0  & 8.1    \\
\bottomrule
\end{tabular}
\end{table}

With a single virtual user the two pipelines differ modestly: the
sequencer posts 5.7~bids/s versus 3.5 for the direct path, a
$1.6\times$ advantage that reflects the absence of lock-wait overhead
even at zero contention. The more telling result emerges at higher load.
At $C\!=\!10$, the direct pipeline slows to 9.5~bids/s while the
sequencer sustains 92.8~bids/s -- nearly a tenfold gap. The p95
latency diverges even more sharply: 1{,}284~ms versus 48~ms.

At $C\!=\!100$ the disparity is pronounced. Direct processing collapses
to 30~bids/s with a p95 of 5{,}772~ms, and the k6 error threshold on
{\tt bid\_errors} was breached, indicating that many requests timed out
or were rejected. The sequencer, by contrast, delivers 975~bids/s at a
p95 of 8.1~ms -- a $\mathbf{32.5\times}$ throughput gain and a $712\times$
latency reduction. The underlying reason is straightforward: a
single-consumer architecture transforms $C$ concurrent lock-waiters
into $C$ entries in a FIFO stream, so the database lock is never
contested regardless of external concurrency.

\subsection{Determinism}

We ran each implementation five times on the same input and computed
SHA-256 over the resulting bid log, ordered by price ascending and row
identifier ascending. Table~\ref{tab:determinism} records how many
distinct hashes appeared.

\begin{table}[h]
\centering
\caption{Bid-Log Determinism Across Five Repeated Runs}
\label{tab:determinism}
\begin{tabular}{lc}
\toprule
\textbf{Implementation} & \textbf{Distinct Hashes} \\
\midrule
Single-Jump          & 1 \\
Recursive            & 3--5 (varies with partial commits) \\
Atomic Ladder        & 1 \\
\bottomrule
\end{tabular}
\end{table}

Both the single-jump and atomic ladder implementations are deterministic
by construction: the former writes one row, the latter writes all rungs
inside a single transaction. The recursive variant is not, since a
partial-commit failure at an arbitrary rung leaves the auction in a
different state from run to run.

\subsection{Log Fidelity}

Table~\ref{tab:fidelity} reports the number of bid rows written per
manual bid when the auto-bid pool should produce a ladder of length
$K^\star = 8$.

\begin{table}[h]
\centering
\caption{Bid Rows per Manual Bid ($K^\star = 8$)}
\label{tab:fidelity}
\begin{tabular}{lrr}
\toprule
\textbf{Implementation} & \textbf{Rows} & \textbf{Fidelity} \\
\midrule
Single-Jump          & 1   & 12.5\% \\
Recursive            & 4--8 (avg 6.3) & 78.8\% \\
Atomic Ladder        & 8   & 100.0\% \\
\bottomrule
\end{tabular}
\end{table}

The recursive approach loses rungs to its partial-commit failure mode: an
abort at rung~4 leaves four rows in the database and an inconsistent
final price. Only the atomic ladder achieves full fidelity on every run.

\subsection{Reproducibility}

All load-test scripts and seed data reside in the project
repository at {\small\tt packages/simulator/k6/\allowbreak
bid-throughput.js}.
Executing {\tt npm run sim:ladder} reproduces every
measurement reported above.

% ─── VII. DISCUSSION ──────────────────────────────────────────────────────────
\section{Discussion}

\subsection{When Single-Jump Is Preferable}

The atomic ladder writes $K$ database rows where single-jump writes one.
For platforms that treat the bid log as an internal ledger rather than a
user-facing narrative -- high-frequency financial exchanges come to
mind -- the extra writes are pure overhead. The protocol is designed for
consumer auction platforms whose product surface includes a live,
scrollable bid history, and its value should be judged in that context.

\subsection{Anti-Sniping Interaction}

A proxy ladder of ten rungs completes in under a millisecond of wall
time, yet the manual bid that triggered it may have landed inside the
auction's anti-sniping window. Our implementation extends the auction
deadline once per manual bid, not once per rung. This design reflects the
user's mental model: the platform is reacting to a single human action,
not to ten separate events.

\subsection{Idempotency Under Job Re-Delivery}

BullMQ's at-least-once delivery guarantee means the ladder job may
execute twice. On re-delivery, the protocol finds that the auction price
has already advanced past every challenger's ceiling and returns an empty
step list. The result is not byte-identical to the first execution --
any manual bid that arrived in the interim will have changed the starting
conditions -- but it is operationally correct: no duplicate rows are
written, and the auction price is never regressed.

\subsection{Threats to Validity}

All measurements were taken on a single host. In a distributed
deployment with a PostgreSQL primary-replica pair and the Socket.io Redis
adapter, network round-trip latencies would increase absolute numbers.
Whether this shifts the {\em relative} ordering of the two pipelines is
an open empirical question.

We also did not test deterministic adversarial worker crashes (e.g., a
forced abort at a specific rung). Single-transaction semantics guarantee
that any such crash leaves no partial state, so the recoverability
argument follows from the atomicity of {\tt COMMIT} rather than from any
property of the protocol itself.

% ─── VIII. CONCLUSION ──────────────────────────────────────────────────────────
\section{Conclusion}

We described the Atomic Ladder Protocol for resolving concurrent proxy
bids in an English auction while preserving a complete per-increment bid
log. The protocol runs every rung of the ladder inside a single database
transaction, obeys a two-tier lock order that eliminates deadlocks, and
produces a final price that matches the Vickrey equilibrium regardless
of bidder arrival order. Five formal properties were stated; three were
proved.

Benchmarks on a PostgreSQL-and-Redis stack indicate that the Redis
Stream sequencer variant sustains 975~bids/s at 100 concurrent bidders,
a $\mathbf{32.5\times}$ improvement over the direct row-locking
baseline, with a 95th-percentile latency of 8.1~ms and 100\% log
fidelity on every run. The protocol is idempotent under worker
re-delivery and has been running in the \textsc{AuctionHaus} platform
since deployment.

Three avenues remain open: (i)~a cross-auction portfolio resolver that
apportions a bidder's available balance across multiple simultaneous
proxy commitments; (ii)~machine-checked verification of the protocol in
TLA$^+$~\cite{Lamport2002}; and (iii)~empirical evaluation under a
multi-node distributed deployment.

% ─── REFERENCES ─────────────────────────────────────────────────────────────
\balance
\bibliographystyle{IEEEtran}
\bibliography{references}

\end{document}
